\question[6]
Il programma \lstinline|guessmyname| (es.\ 09) contiene \lstinline|static const char *secret_answer = "42";| e confronta l'input con \lstinline|strcmp|.
(a) In quale sezione ELF finisce la stringa \lstinline|"42"|? (b) Perché \lstinline|strings guessmyname| di default non la mostra? (c) Indica due comandi che rivelano il segreto.
\begin{solution}
(a) In \lstinline|.rodata| (costante di sola lettura).\\
(b) \lstinline|strings| di default stampa solo sequenze di almeno 4 caratteri; \lstinline|"42"| ne ha 2.\\
(c) \lstinline|readelf -x .rodata guessmyname| (nel dump compare \lstinline|3432| = ``42'') oppure \lstinline|strings -n 2 guessmyname|. Verifica: \lstinline!echo 42 | ./guessmyname! $\to$ ``Esatto!''.
\end{solution}
